% demo.tex — exercises every template of the AINIT Academic Beamer theme.
% Compile with: lualatex demo.tex (run twice for cross-references).

\documentclass[aspectratio=169,11pt]{beamer}

% Load our theme. The .sty files live in the same directory as this demo.
\usetheme{AINITAcademic}

\usepackage{appendixnumberbeamer}

% Lecture identity
\title{Two Engines}
\subtitle{Memory, prediction, and the difference between them}
\author{\textsc{Carlos Cotrini}}
\institute{ETH Zürich \,|\, CAS AINIT FS26}
\date{Week 2 \textemdash{} Spring 2026}

% Page-furniture defaults; per-frame overrides below.
\renewcommand{\AINITshorttitle}{Two engines}
\renewcommand{\AINITsection}{}

\begin{document}

% =============================== 1. TITLE ================================
\AINITmaketitle

% =========================== 2. SECTION DIVIDER ==========================
\sectiondivider%
  {``A man's memory is bound to be a distraction unless it can be discarded;\\ then a sleeper or a savant might think clearly.''}%
  {Jorge Luis Borges, \textit{Funes the Memorious}}%
  {Act I \textemdash{} Memory}

% =============================== 3. CONTENT ==============================
\renewcommand{\AINITsection}{Act I \textemdash{} Memory}
\begin{frame}{Two engines, one question}
  \sidenote{The distinction is older than statistics: see Hume's \textit{Treatise} on causation versus constant conjunction.}
  We will spend the next ninety minutes inside a single question:

  \begin{itemize}
    \item what does it mean for a system to \AINITemph{remember}, as opposed to \AINITemph{predict}?
    \item where in your organisation is each engine appropriate?
    \item and what changes when the data is sparse, drifting, or adversarial?
  \end{itemize}

  The two engines look the same from the outside. They are not the same on the inside, and the difference will matter to every procurement decision you make this year.
\end{frame}

% ============================== 4. EQUATION ==============================
\renewcommand{\AINITsection}{Act II \textemdash{} Prediction}
\equationslide%
  {Cross-entropy, decomposed}%
  {\text{cross-entropy} \;=\; H(p) \;+\; {\color{accentOxblood}\mathrm{KL}\!\left(p \,\|\, \hat p\right)}}%
  {Only the second term is yours to minimise; the first is the world's irreducible surprise.}

% ============================ 5. VIZ HANDOFF =============================
\vizhandoff{Customer-urns carousel}

% ============================== 6. CLOSING ===============================
% Section indicator hidden on the closing slide for cleanliness.
\renewcommand{\AINITsection}{}
\closingslide%
  {Funes remembered everything and could think nothing.}%
  {Machine learning forgets almost everything \textemdash{} and that is why it understands.}

% ============================ 7. REFERENCES ==============================
\appendix
\renewcommand{\AINITshorttitle}{Two engines \textemdash{} references}
\begin{frame}{References}
  \footnotesize
  \begin{enumerate}
    \item J.~L. Borges, ``Funes el memorioso,'' in \textit{Ficciones}, 1944.
    \item D.~Hume, \textit{A Treatise of Human Nature}, Book I, 1739.
    \item T.~M. Cover and J.~A. Thomas, \textit{Elements of Information Theory}, 2nd ed., Wiley, 2006.
    \item C.~M. Bishop, \textit{Pattern Recognition and Machine Learning}, Springer, 2006, ch.~1.
    \item E.~T. Jaynes, \textit{Probability Theory: The Logic of Science}, CUP, 2003.
  \end{enumerate}
\end{frame}

\end{document}
